\section{Queues}
\label{sec:alps-queues}

Durch die Integration von \gls{alps} verlängert sich die \textit{End-to-End-Verarbeitungszeit} einer Alarmmeldung, also die Zeit vom Senden der Alarmmeldung durch die Alarmanlage bis zum Ankommen bei \gls{srvt}. Dafür gibt es zwei Gründe. Erstens werden der Verarbeitungskette vier zusätzliche Komponenten hinzugefügt: \gls{alpone}, der Broker, \gls{alptwo} und \gls{alpdedup}. Zweitens wirken sich \texttt{x-single-active-\allowbreak consumer} und \texttt{prefetch\_count = 1} negativ auf die Verarbeitungszeit aus, da pro Queue immer nur eine Alarmmeldung gleichzeitig verarbeitet wird. Die Anzahl an Alarmmeldungen, die parallel verarbeitet werden können, hängt somit von der Anzahl der Queues ab. Alarmmeldungen verschiedener Alarmanlagen, die in derselben Queue landen, müssen aufeinander warten, obwohl zwischen ihnen keine Reihenfolge eingehalten werden muss. Um diese Verschlechterung zu kompensieren, bietet sich die Möglichkeit, die Anzahl der Queues zu erhöhen.

Der \texttt{x-consistent-hash} Exchange verwendet den Algorithmus \textit{Jump Consistent Hash}~\citep{lamping2014}. \gls{alpone} nutzt die \texttt{device\_id} als \texttt{routing-key}, also als den Wert, der gehasht wird. Der Algorithmus teilt den Hash-Raum nahezu perfekt gleichmäßig auf die Queues auf~\citep{lamping2014}. Daraus lässt sich jedoch nicht ableiten, welche Anzahl an Queues ideal für \gls{alps} ist. Die Gründe dafür sind:

\begin{itemize}
    \item \textbf{Endliche Anzahl an Alarmanlagen}: Die gleichmäßige Aufteilung bezieht sich auf den gesamten Hash-Raum. In \gls{alps} wird jedoch nur eine endliche Anzahl an \texttt{device\_id}s verteilt. Die \texttt{device\_id}s sind sechsstellige Ziffernfolgen. Wie gleichmäßig genau diese \texttt{device\_id}s auf die Queues verteilt werden, ist nicht durch den Algorithmus festgelegt. Je mehr Queues es gibt, desto weniger Alarmanlagen bekommt jede Queue und desto größer werden die relativen Unterschiede zwischen den Queues.

    \item \textbf{Ungleichmäßige Last}: Eine gleichmäßige Verteilung der Alarmanlagen bedeutet keine gleichmäßige Verteilung der Alarmmeldungen. Alarmanlagen senden unterschiedlich viele Alarmmeldungen zu unterschiedlichen Zeitpunkten. Landen zwei gleichzeitig aktive Alarmanlagen in derselben Queue, entsteht eine Wartezeit, auch wenn die Verteilung insgesamt gleichmäßig ist.

    \item \textbf{Zusätzlicher Aufwand}: Jede zusätzliche Queue erhöht die Anzahl der Consumer und der Verbindungen zum Broker. Gleichzeitig erzeugen mehr parallele Alarmmeldungen mehr parallele Requests an das Legacy-System, dessen Verarbeitungszeit laut dem Vorhersagemodell von der Anzahl paralleler Requests abhängt (siehe Abschnitt~\ref{sub-sub:exp-alps-phase-3}). Ab einer bestimmten Anzahl an Queues kann sich die Verarbeitungszeit daher nicht mehr verbessern oder sogar verschlechtern.
\end{itemize}

Aus diesen Gründen wird die ideale Anzahl an Queues experimentell bestimmt. Dafür wird das Experiment aus \textit{Phase~3} (Abschnitt~\ref{sub-sub:exp-alps-phase-3}) mit steigender Anzahl an Queues $N_{\text{queue}}$ wiederholt, während alle übrigen Parameter unverändert bleiben. Für jede Anzahl an Queues wird die \textit{End-to-End-Verarbeitungszeit} der Alarmmeldungen gemessen und mit der von \gls{lps} verglichen.
